\question[4]
Gegeben ist die Zustandsübergangstabelle eines Mealy-Automaten mit Eingabe- und Ausgabealphabet $\{0,1\}$.
In jeder Zelle steht \enquote{Folgezustand / Ausgabe}.

\begin{center}
\begin{tabular}{r|cc}
 & 0 & 1 \\ \hline
$\rightarrow s_0$ & $s_0$ / 0 & $s_1$ / 1 \\
$s_1$ & $s_0$ / 1 & $s_1$ / 0 \\
\end{tabular}
\end{center}

\begin{parts}
\part Zeichne das Zustandsdiagramm.
\part Bestimme die Ausgabe für die Eingabe $0011101$.
\part Beschreibe, wann der Automat eine 1 ausgibt.
\end{parts}

\begin{solutionbox}{6cm}
\begin{parts}
\part \mbox{}\\
\includegraphics[height=2.5cm]{\pfad/mealy_02/mealy_02.png}
\part $0010011$
\part Der Automat gibt eine 1 aus, wenn sich das Eingabesymbol gegenüber dem vorherigen ändert
(am Anfang wird eine 0 als vorheriges Symbol angenommen). Im Zustand $s_0$ war das letzte Symbol eine 0,
im Zustand $s_1$ eine 1.
\end{parts}
\end{solutionbox}
